\documentclass[11pt]{article}
\usepackage[a4paper,margin=0.9in]{geometry}
\usepackage[T1]{fontenc}
\usepackage{helvet}
\renewcommand{\familydefault}{\sfdefault}
\usepackage{xcolor}
\usepackage{booktabs}
\usepackage{array}
\usepackage{colortbl}
\usepackage{float}
\usepackage{amsmath}
\usepackage{titlesec}
\usepackage{fancyhdr}
\usepackage{enumitem}
\usepackage{lastpage}
\usepackage[hidelinks]{hyperref}

\definecolor{bcnavy}{RGB}{11,31,58}
\definecolor{bcblue}{RGB}{32,74,135}
\definecolor{bcgold}{RGB}{176,141,87}
\definecolor{bcgrey}{RGB}{90,98,112}
\definecolor{bcrule}{RGB}{210,214,220}
\definecolor{bckeep}{RGB}{232,237,244}
\definecolor{bcact}{RGB}{247,241,228}
\definecolor{bcwarn}{RGB}{250,238,238}

\titleformat{\section}{\color{bcnavy}\Large\bfseries}{\thesection}{0.6em}{}
\titleformat{\subsection}{\color{bcblue}\large\bfseries}{\thesubsection}{0.6em}{}
\titlespacing*{\section}{0pt}{1.3em}{0.45em}

\pagestyle{fancy}\fancyhf{}
\renewcommand{\headrulewidth}{0pt}
\renewcommand{\footrulewidth}{0.4pt}
\renewcommand{\footrule}{\hbox to\headwidth{\color{bcrule}\leaders\hrule height \footrulewidth\hfill}}
\fancyfoot[L]{\footnotesize\color{bcgrey}Bayesian Capital}
\fancyfoot[C]{\footnotesize\color{bcgrey}Confidential --- Internal Research}
\fancyfoot[R]{\footnotesize\color{bcgrey}Page \thepage\ of \pageref{LastPage}}

\setlength{\parindent}{0pt}\setlength{\parskip}{0.5em}

\newcommand{\note}[2]{%
  \begin{center}\fcolorbox{bcrule}{#1}{\parbox{0.95\textwidth}{\vspace{0.2em}#2\vspace{0.2em}}}\end{center}}

\begin{document}

{\color{bcgold}\rule{\textwidth}{1.6pt}}\\[0.3cm]
{\color{bcnavy}\huge\bfseries Live Daily Book}\\[0.15cm]
{\color{bcblue}\Large Parameters, Allocation and Return Expectation}\\[0.35cm]
{\color{bcgold}\rule{\textwidth}{0.8pt}}\\[0.35cm]

\begin{tabular}{@{}ll@{}}
\textbf{Workbook} & \texttt{TradingExcel\_5stock\_live.xlsx} \\
\textbf{Date of this note} & 5 August 2026 \\
\textbf{Price history in the file} & 1 April 2024 to 3 August 2026 --- 587 sessions, 2.34 years \\
\textbf{Names} & TSM, VRT, VST, RKLB, MU (daily model) \\
\textbf{Not funded from this workbook} & NVDA and AVGO, which run the weekly Monday model \\
\end{tabular}

\vspace{0.6em}
{\color{bcrule}\rule{\textwidth}{0.6pt}}

\section{What the workbook is}

Five names, each traded by two tranches. The Bayes tranche bids below a Kalman
local-linear-trend estimate of fair value; the OU tranche bids below a one-step AR(1) forecast
taken over a rolling window. Both bids are then floored by the day's open and by a fixed
percentage below the running all-time high, and both sell at the bid plus a premium on the prior
close. A position that has not reached its target within 50 calendar days is closed at the open.
Cash earns 3.14\% while it waits and commission is half a cent a share.

The two tranches share the name's capital in fixed proportion. That proportion is now 50/50 on
every name: the 75\% tilt towards Bayes was compensating for a mis-specified OU sigma, and once
the sigma was corrected the tilt won only 7 of 15 walk-forward folds and lost 2.1pp.

\section{Parameter set}

Everything below is read from the \texttt{Model} sheet of each name. Cell references are given so
the table can be checked against the file directly.

\vspace{0.4em}
\begin{center}
\footnotesize
\begin{tabular}{@{}l r r r r r r r r r r@{}}
\toprule
& \multicolumn{6}{c}{\textbf{Bayes tranche}} & \multicolumn{4}{c}{\textbf{OU tranche}} \\
\cmidrule(lr){2-7}\cmidrule(lr){8-11}
& $\lambda$ & $\varphi_L$ & $\psi$ & $k$ & prem. & peak cap & $W$ & buffer & prem. & cap \\
& \tiny B2 & \tiny D2 & \tiny F2 & \tiny H2 & \tiny J2 & \tiny L2 & \tiny B3 & \tiny D3 & \tiny H3 & \tiny J3 \\
\midrule
TSM  & 0.4620 & 0.2543 & 0.0650 & 1.0045 & 0.01565 & 0.03048 & 77  & 0.20 & 0.01002 & 0.02089 \\
VRT  & 0.4207 & 0.4259 & 0.0817 & 1.1164 & 0.02169 & 0.04838 & 80  & 0.40 & 0.02660 & 0.05950 \\
VST  & 0.3735 & 0.2643 & 0.0962 & 1.3857 & 0.02182 & 0.01366 & 122 & 0.65 & 0.01450 & 0.05275 \\
RKLB & 0.3528 & 0.2374 & 0.1130 & 0.7123 & 0.02684 & 0.00802 & 85  & 0.25 & 0.02215 & 0.03365 \\
MU   & 0.6000 & 0.2636 & 0.0100 & 1.0548 & 0.02430 & 0.03800 & 91  & 0.75 & 0.02500 & 0.02920 \\
\bottomrule
\end{tabular}
\end{center}

Shared across all five names: commission \$0.005 per share (N2), IBKR interest 3.14\% p.a.\ (R2),
stop-loss 50 calendar days (T2), Bayes share 0.50 (V2). Each \texttt{Model} sheet backtests on a
notional \$1{,}000{,}000 (P2); the live capital comes from the \texttt{Allocation} sheet and P2
does not need changing.

\subsection{The OU buffer is on the residual scale}

The buffers in D3 look inconsistent with earlier versions of the book because the quantity they
multiply changed. Column AZ now measures the standard deviation of the fitted AR(1)
\emph{residuals}, not \texttt{STDEVP} of the last $W$ closes. The old measure was the dispersion
of the price level, which on a trending name is mostly the trend --- so the buffer widened
exactly when a stock was running and fills were lost. Residual sigma is roughly a third of the
old value, and D3 was raised to compensate. The pairs are not interchangeable: a buffer from one
scale used with a sigma from the other gives a bid that is wrong by a factor of about three.
Section~6 is about a place in the workbook where exactly that has happened.

\section{Allocation settings}

\begin{center}
\footnotesize
\begin{tabular}{@{}l l r l@{}}
\toprule
\textbf{Cell} & \textbf{Input} & \textbf{Value} & \textbf{Effect} \\
\midrule
B5 & Total capital available & \$2{,}782{,}455.86 & Split across the five daily names only \\
B6 & Bayes fraction, flat    & 0.50 & \emph{Not read} while B9 = 2 --- see below \\
B7 & Volatility lookback     & 120 days & Feeds the average daily range in column D \\
B8 & Floor, min weight       & 0.20 & \\
E11:E15 & Cap, max weight    & 0.20 & Floor = cap, so every weight is pinned to 0.20 \\
B9 & Split mode              & 2 & Per-stock; reads Q11:Q15, all of which are 0.50 \\
C11:C15 & Return assumption  & 0.55 / 0.53 / 0.57 / 1.15 / 0.54 & Inert while floor = cap \\
\bottomrule
\end{tabular}
\end{center}

The floor and the cap are both 0.20, so the risk-weighting machinery in columns F to K runs but
cannot change anything: each name gets exactly one fifth. That is deliberate --- equal weight was
the decision --- but it means the return assumptions in column C and the volatility lookback in
B7 have no effect on the dollar allocation at all. If the weighting is ever re-enabled by
loosening E11:E15, column C becomes live and should be refreshed to the planning figures in
Section~4 first; the values in there now are stale.

\vspace{0.3em}
Resulting allocation, per name:

\begin{center}
\footnotesize
\begin{tabular}{@{}l r r r@{}}
\toprule
& \textbf{Weight} & \textbf{\$ to the name} & \textbf{\$ per tranche} \\
\midrule
TSM, VRT, VST, RKLB, MU & 0.20 each & 556{,}491.17 & 278{,}245.59 Bayes / 278{,}245.59 OU \\
\midrule
\textbf{Total} & 1.00 & \textbf{2{,}782{,}455.86} & 10 tranches \\
\bottomrule
\end{tabular}
\end{center}

\note{bcact}{\textbf{One trap worth knowing about.} Split mode (B9) is set to 2, so the Bayes
share actually used is column Q, not cell B6. Q11:Q15 are all 0.50 and B6 is also 0.50, so the
book behaves correctly today --- but editing B6 will change nothing. To move the split, either
set B9 to 1 or edit Q11:Q15, and mirror the change into V2 on each \texttt{Model} sheet, which is
what the backtest reads.}

\section{What return to expect}

\subsection{The planning figures}

These are annualised, verified against intraday bars, marked to market, at the deployed 50/50
split with the corrected OU sleeve.

\begin{center}
\footnotesize
\begin{tabular}{@{}l r r r r@{}}
\toprule
& \textbf{First half} & \textbf{Tested half} & \textbf{Full sample} & \textbf{Planning figure} \\
& \tiny to 23 May 2025 & \tiny from 23 May 2025 & & \tiny full sample less the tilt haircut \\
\midrule
RKLB & 162\% & 154\% & 158\% & 156\% \\
TSM  &  34\% &  82\% &  57\% &  55\% \\
VST  &  53\% &  71\% &  62\% &  60\% \\
VRT  & --5\% & 192\% &  67\% &  65\% \\
MU   &   7\% & 142\% &  63\% &  61\% \\
\midrule
\textbf{Book, equal weight} & \textbf{50\%} & \textbf{128\%} & \textbf{81\%} & \textbf{79\%} \\
\textbf{Book excluding RKLB} & \textbf{22\%} & \textbf{122\%} & \textbf{62\%} & \textbf{60\%} \\
\bottomrule
\end{tabular}
\end{center}

\note{bckeep}{\textbf{The number to plan on is 79\% a year, and it should be carried with a wide
band --- call it 60\% to 80\%.} Two things pull it down. RKLB alone contributes 31 of the 79
points, so the book's headline depends on one name that has returned 156\% and is not going to be
assumed to do so indefinitely; without it the four remaining names average 60\%. And the
dispersion between halves is large: 50\% in the first half against 128\% in the tested half. That
the tested half is the better one is reassuring about the fit, but it is not evidence the higher
level persists. The full-sample figure is the middle course between the two regimes and that is
why it is the planning basis.}

On \$2{,}782{,}455.86 the 79\% central figure is about \$2.2m a year and the 60\% floor about
\$1.7m. Both are before tax and before any slippage beyond the modelled half-cent commission.

\subsection{Do not use the workbook's own Annual return cell}

Cell Y5 on each \texttt{Model} sheet reports 143\% for TSM and 508\% for RKLB. Those figures are
arithmetically right and are not a forecast. The sheet books a completed round trip whenever the
day's low reached the bid \emph{and} the day's high reached the target, because with daily bars it
has no way to know the low came first. On these names, trading a 3\% daily range against a
premium of 1.6\% to 2.7\%, that assumption is doing a great deal of work.

Re-running the identical model in Python under each convention isolates it:

\begin{center}
\footnotesize
\begin{tabular}{@{}l r r r r@{}}
\toprule
& \textbf{Sheet cell Y5} & \textbf{Sheet's rule} & \textbf{Verified fills} & \textbf{No same-day exits} \\
& & \tiny all same-day exits & \tiny 5-minute bars & \tiny most conservative \\
\midrule
RKLB & 508\% & 438\% & \textbf{158\%} & 36\% \\
TSM  & 143\% & 129\% & \textbf{~57\%} & 39\% \\
VST  & 263\% & 240\% & \textbf{~62\%} & 23\% \\
VRT  & 242\% & 208\% & \textbf{~67\%} & 35\% \\
MU   & 178\% & 135\% & \textbf{~63\%} & 41\% \\
\bottomrule
\end{tabular}
\end{center}

The gap between the first two columns is not a modelling difference: Y5 annualises with a
hard-coded $^{1/2.2}$ exponent while the feed now spans 2.34 years. Correct that and the Python
replica reproduces the sheet to within a point. The whole of the remaining drop is same-day fill
verification, and the middle column is the honest one.

\subsection{Activity}

\begin{center}
\footnotesize
\begin{tabular}{@{}l r r r r@{}}
\toprule
& \textbf{Buys / yr} & \textbf{Bayes} & \textbf{OU} & \textbf{Stops in 2.3 yrs} \\
\midrule
RKLB & 96 & 52 & 44 & 5 \\
TSM  & 65 & 42 & 23 & 3 \\
VST  & 68 & 37 & 31 & 5 \\
VRT  & 54 & 41 & 13 & 4 \\
MU   & 51 & 28 & 23 & 4 \\
\midrule
\textbf{Book} & \textbf{335} & 200 & 135 & 21 \\
\bottomrule
\end{tabular}
\end{center}

About 335 buys and a similar number of sells a year: 1.3 buys on an average trading day. A
tranche in cash bids every morning and the tranches are in cash 56\% of the time, so roughly 5.6
of the 10 orders go in on a typical day. The 50-day stop fires about nine times a year across the
book.

\section{Daily operation}

Each morning, read the \texttt{Dashboard}: for any tranche currently in cash, place a limit buy at
its buy level and a GTC sell at its target. Tranches already holding stock need nothing. The
\texttt{Allocation} sheet's machine-readable block, rows 25 to 34, gives the dollar amount per
stock and tranche for the IBKR script.

\section{Two things to correct before this file is traded}

\note{bcwarn}{\textbf{1. The Dashboard's OU sigma was never moved to the residual definition,
and it is the sheet the orders are read from.} Column P still computes
\texttt{STDEVP(last $W$ closes)} --- the level dispersion Section~2.1 describes --- while the
\texttt{Model} sheets use residual sigma and the buffers in D3 were raised on the assumption they
would be multiplying the smaller quantity. The two errors compound, so the OU buy levels come out
far below the market and would not fill:

\vspace{0.3em}
\begin{center}
\footnotesize
\begin{tabular}{@{}l r r r r r@{}}
\toprule
& \textbf{Close} & \textbf{Sigma shown} & \textbf{Correct sigma} & \textbf{OU buy shown} & \textbf{OU buy correct} \\
\midrule
TSM  & 406.11 & 23.63 & 12.42 & 402.73 & 404.97 \\
VRT  & 263.05 & 27.41 & 13.69 & 258.05 & 263.54 \\
VST  & 155.94 &  8.32 &  4.49 & 150.79 & 153.28 \\
RKLB &  70.43 & 23.76 &  6.54 &  65.42 &  69.72 \\
MU   & 829.50 & 254.99 & 56.20 & \textbf{635.96} & \textbf{785.06} \\
\bottomrule
\end{tabular}
\end{center}
\vspace{0.2em}
MU's OU tranche has been bidding 23\% below the market. The backtests are unaffected --- the
\texttt{Model} sheets are correct --- so nothing in Section~4 changes. A corrected copy is
supplied as \texttt{TradingExcel\_5stock\_live\_fixed.xlsx}, in which Dashboard P4:P8 use the same
residual construction as \texttt{Model!AZ}.}

\note{bcact}{\textbf{2. The \texttt{Notes} and \texttt{Allocation} text still say 0.75.}
\texttt{Notes} B4 and \texttt{Allocation} A2 describe a Bayes share of 0.75 in flat mode. The
cells are set to 0.50 in per-stock mode and the book behaves correctly; only the prose is stale.
It is worth fixing before anyone reads the sheet as documentation.}

\vspace{0.8em}
{\color{bcrule}\rule{\textwidth}{0.6pt}}

{\footnotesize\color{bcgrey}
All figures verified against 5-minute bars, marked to market, net of \$0.005 per share commission
and inclusive of interest on idle cash at 3.14\%. Past backtested performance is not a forecast of
future returns; the sample is 2.3 years and covers one broadly rising market.}

\end{document}
